% !TeX program = pdflatex
\documentclass[12pt,a4paper]{article}

\newcommand{\notelanguage}{finnish}
% Shared preamble for course notes and exercise sets.

\usepackage[T1]{fontenc}
\usepackage{lmodern}
\usepackage[english]{babel}

\usepackage[a4paper,margin=30mm]{geometry}
\usepackage{microtype}
\usepackage{graphicx}
\usepackage{enumitem}

\usepackage{amsmath}
\usepackage{amssymb}
\usepackage{amsthm}
\usepackage{mathtools}
\usepackage{aliascnt}

\usepackage[hidelinks,pdfusetitle]{hyperref}

\numberwithin{equation}{section}
\setlist{itemsep=0.25em,topsep=0.5em}

\theoremstyle{plain}
\newtheorem{theorem}{Theorem}[section]
\newaliascnt{lemma}{theorem}
\newtheorem{lemma}[lemma]{Lemma}
\aliascntresetthe{lemma}
\newaliascnt{proposition}{theorem}
\newtheorem{proposition}[proposition]{Proposition}
\aliascntresetthe{proposition}
\newaliascnt{corollary}{theorem}
\newtheorem{corollary}[corollary]{Corollary}
\aliascntresetthe{corollary}

\theoremstyle{definition}
\newaliascnt{definition}{theorem}
\newtheorem{definition}[definition]{Definition}
\aliascntresetthe{definition}
\newaliascnt{example}{theorem}
\newtheorem{example}[example]{Example}
\aliascntresetthe{example}
\newaliascnt{problem}{theorem}
\newtheorem{problem}[problem]{Problem}
\aliascntresetthe{problem}

\theoremstyle{remark}
\newaliascnt{remark}{theorem}
\newtheorem{remark}[remark]{Remark}
\aliascntresetthe{remark}

\usepackage[nameinlink,noabbrev]{cleveref}

\crefname{theorem}{theorem}{theorems}
\crefname{lemma}{lemma}{lemmas}
\crefname{proposition}{proposition}{propositions}
\crefname{corollary}{corollary}{corollaries}
\crefname{definition}{definition}{definitions}
\crefname{example}{example}{examples}
\crefname{problem}{problem}{problems}
\crefname{remark}{remark}{remarks}

\DeclareMathOperator{\im}{im}
\DeclareMathOperator{\rank}{rank}
\DeclarePairedDelimiter{\abs}{\lvert}{\rvert}
\DeclarePairedDelimiter{\norm}{\lVert}{\rVert}
\DeclarePairedDelimiter{\set}{\lbrace}{\rbrace}


\title{Harjoituksen 1 ratkaisut}
\author{}
\date{}

\begin{document}

\maketitle

\noindent

\begin{exercise}
  Merkitään propositiosymboleilla \(p_0\), \(p_1\) ja \(p_2\) seuraavia
  lauseita:
  \begin{align*}
    p_0&: \text{Sataa.} \\
    p_1&: \text{Tuulee.} \\
    p_2&: \text{On Kylmä.}
  \end{align*}
  Esitä luonnollisella kielellä lause
  \[\left((p_0 \land \neg p_1) \rightarrow \neg p_2\right).\]
\end{exercise}

\begin{solution}
\end{solution}

\begin{exercise}
  Annetaan propositiosymboleille merkitykset tehtävän 1 tapaan. Käännä
  seuraavat lauseet propositiologiikan kielelle:
\begin{enumerate}[label=\alph*), nosep]
  \item Jos on kylmä, mutta ei sada, niin tuulee.
  \item Jos ei sada, niin ei ole kylmä, paitsi jos tuulee.
\end{enumerate}
\end{exercise}

\begin{solution}
\end{solution}

\begin{exercise}
Tarkastellaan lausetta ``Jos Niilo Syö valtavasti aina, kun Unicafessa
on linssibolognesea, niin joko Niilo on kova syömään tai linssibolognese on
hänen lempiruokansa.``
\begin{enumerate}[label=\alph*), nosep]
\item Mitkä ovat lauseen sisältämät atomilauseet?
\item Formalisoi lause propositiolauseella. (Huomaa, että joudut korvaamaan
    luonnollisen kielen vivahteikkaampia ilmaisuja propositiologiikan
  konnektiiveilla.)
\end{enumerate}
\end{exercise}

\begin{solution}
\end{solution}

\begin{exercise}
Keksi luonnollisen kielen lause, jonka formalisointi on muotoa
\[(((p_0 \rightarrow (p_1 \lor p_2)) \land \neg p_2) \rightarrow (p_0
\rightarrow p_1)).\]
\end{exercise}

\begin{solution}
\end{solution}

\begin{exercise}
Mitkä seuraavista merkkijonoista ovat propositiolauseita? Perustele.
\begin{enumerate}[label=\alph*), nosep]
\item \(((p_0 \land p_0) \land \neg p_0)\)
\item \(p_0 \rightarrow p_1 \rightarrow p_2\)
\item \(p_{32110}\)
\item \((p_0) \lor (\neg p_0)\)
\item \((\neg(\neg p_{222}))\)
\item \())p_1 \lor \rightarrow p_0\)
\item \((p_5 \rightarrow (p_0 \rightarrow \neg \neg p_0))\)
\item \((p_0 \neg \land p_2)\)
\end{enumerate}
\end{exercise}

\begin{solution}
\end{solution}

\end{document}
